\section{}

\begin{frame}[t]{講義・実習の目的}
  \begin{itemize}
    %\setlength{\itemsep}{1em}
  \item 理論・実験を問わず、学部〜大学院〜で必要となる現代的かつ普遍的な計算機の素養を身につける
  \item より高度な数値計算アルゴリズム・数値計算の常識を学ぶ
  \item 物理学における具体的な問題を通して実践的な知識と経験を身につける\\[2em]
  \item 数学公式と数値計算アルゴリズムは別物であることを理解する
  \item 刻み幅・近似度合い・サンプル数・初期状態・乱数seed・ビン幅を変えて結果を確認する
  \item 計算量(コスト)のスケーリング(次数)に気をつける
  \item いろいろな手法を組み合わせて問題を解決する
  \end{itemize}
\end{frame}

\begin{frame}[t]{講義・実習内容}
  \begin{itemize}
    \setlength{\itemsep}{1em}
  \item 問題解決型: 計算機実験Iで身に付けた知識をもとに、より高度な数値計算手法・アルゴリズムを学び、物理学における具体的な問題への応用を通して実践的な知識と経験を身につける
    \begin{itemize}
    \item 取り上げる問題
      \begin{itemize}
      \item ランダムプロセス
      \item 統計力学(イジング模型)
      \item 拡散方程式、波動方程式
      \item 量子力学(1粒子系、実時間発展、横磁場イジング模型、量子コンピュータ)
      \item 力学系・粒子系
        
      \end{itemize}
    \item 取り上げる計算手法・アルゴリズム
      \begin{itemize}
      \item モンテカルロ法・マルコフ連鎖モンテカルロ法
      \item 行列の対角化・疎行列分解
      \item 有限差分法、偏微分方程式の解法
      \item 常微分方程式の積分、分子動力学
      \item 連続最適化、離散最適化
      \end{itemize}
    \end{itemize}
    などを予定
  \end{itemize}
\end{frame}

\begin{frame}[t,fragile]{講義と実習}
  \begin{itemize}
    %\setlength{\itemsep}{1em}
  \item スタッフ \href{mailto:computer@phys.s.u-tokyo.ac.jp}{computer@phys.s.u-tokyo.ac.jp}
    \begin{itemize}
    \item 講義: 藤堂
    \item 実習: 諏訪助教(藤堂研)、高橋助教(樺島研)
    \item 実習TA: 所(藤堂研)
    \end{itemize}
  \item 講義・実習の進め方
    \begin{itemize}
    \item 講義(座学)と実習を実施
    \item 実習の回には各自PCを持参すること
    \end{itemize}
  \item 評価
    \begin{itemize}
    \item レポート(計2回)
    \item グループ内プレゼンテーション(実習3,4の回に実施予定)
    \item 出席（実習の回）
    \end{itemize}    
  \end{itemize}
\end{frame}
